% ============================================================================
% LEHRERHINWEISE: Kräfte zusammensetzen (UE 2)
% Physik Klasse 7 – Gesamtschule MV – 45 Minuten
% ============================================================================
\documentclass[a4paper,11pt]{article}

\usepackage[left=2cm,right=2cm,top=2cm,bottom=2cm]{geometry}
\usepackage[utf8]{inputenc}
\usepackage[T1]{fontenc}
\usepackage[ngerman]{babel}
\usepackage{array}
\usepackage{tabularx}
\usepackage{fancyhdr}
\usepackage{tcolorbox}
\usepackage{enumitem}
\usepackage{xcolor}

\renewcommand{\familydefault}{\sfdefault}

\pagestyle{fancy}
\fancyhf{}
\renewcommand{\headrulewidth}{0.4pt}
\lhead{\textbf{LEHRERHINWEISE}}
\chead{Kräfte zusammensetzen}
\rhead{Physik Klasse 7}
\setlength{\headheight}{14pt}
\cfoot{\thepage}

\definecolor{physik}{HTML}{2c5aa0}
\definecolor{lightgray}{RGB}{245,245,245}

\setlength{\parindent}{0pt}

\begin{document}

\section*{Lehrerhinweise: Kräfte zusammensetzen (UE 2)}

\textbf{Dauer:} 45 Minuten \quad | \quad \textbf{Materialien:} LP, AB, SIM, MT

\vspace{0.5cm}

\subsection*{1. Stundenverlauf}

\begin{tabular}{|c|p{2.5cm}|p{5cm}|p{4.5cm}|}
\hline
\textbf{Min} & \textbf{Phase} & \textbf{Inhalt} & \textbf{Medien/Sozialform} \\
\hline
0--5 & Einstieg & Wiederholung Kraftpfeil: 3 Elemente (Angriffspunkt, Richtung, Betrag) & LP Schritt 1, EA \\
\hline
5--12 & Erarbeitung I & Maßstabszeichnen: 1 cm $\hat{=}$ 1 N, Pfeile zeichnen & LP Schritt 2 + AB Aufg. 1, EA \\
\hline
12--20 & Erarbeitung II & Resultierende bei gleicher Richtung: $F_{res} = F_1 + F_2$ & LP Schritt 3 + SIM, PA \\
\hline
20--28 & Übung & Tauziehen: entgegengesetzte Kräfte, $F_{res} = |F_1 - F_2|$ & LP Schritt 4--5 + AB Aufg. 3, PA \\
\hline
28--35 & Vertiefung & Rechtwinklige Kräfte (Bonus): Pythagoras & LP Schritt 6 + SIM, EA \\
\hline
35--45 & Sicherung & Minitest (10 Min), Abgabe & MT, EA \\
\hline
\end{tabular}

\vspace{0.5cm}

\subsection*{2. Differenzierung}

\begin{tcolorbox}[colback=lightgray,colframe=physik,boxrule=1pt,arc=0pt]
\textbf{Basisniveau ($\star$):} AB Aufgaben 1--4 (15 BE), MT Aufgaben 1--4 (13 BE)

\textbf{Erweiterung ($\star\star$):} Aufgabe 5 Rechtwinklig (+3 BE, alle machen durch LP)

\textbf{Schnelle SuS:} Eigene Szenarien in der SIM ausprobieren
\end{tcolorbox}

\vspace{0.3cm}

\subsection*{3. Häufige Fehler und Reaktionen}

\begin{tabular}{|p{6cm}|p{7.5cm}|}
\hline
\textbf{Typischer Fehler} & \textbf{Reaktion/Hilfestellung} \\
\hline
Pfeillänge stimmt nicht mit Maßstab überein & ``Zähle die Kästchen -- wie viele cm sind das?'' \\
\hline
Entgegengesetzte Kräfte werden addiert & ``Ziehen beide in dieselbe Richtung?'' \\
\hline
Richtung der Resultierenden falsch & ``Wer ist stärker beim Tauziehen?'' \\
\hline
Gleichgewicht ($F_{res}=0$) nicht erkannt & ``Was passiert, wenn beide gleich stark sind?'' \\
\hline
Pythagoras-Formel falsch angewendet & ``Erst quadrieren, dann addieren, dann Wurzel'' \\
\hline
\end{tabular}

\vspace{0.5cm}

\subsection*{4. Material-Checkliste}

\begin{itemize}[nosep]
\item[$\square$] Beamer/Smartboard für SIM (Demonstration)
\item[$\square$] Schüler-Tablets/Laptops für LP und MT
\item[$\square$] AB ausgedruckt (1 pro Schüler)
\item[$\square$] Lineale (für Maßstabszeichnen)
\item[$\square$] QR-Codes für LP und MT bereithalten
\end{itemize}

\vspace{0.5cm}

\subsection*{5. Hinweise zur Simulation (SIM)}

\begin{itemize}[nosep]
\item Die SIM zeigt zwei verschiebbare Kraftvektoren auf einem Rastergitter
\item 1 Kästchen = 1 N (Maßstab wie im AB)
\item Resultierende wird automatisch berechnet und angezeigt
\item \textbf{Voreinstellungen nutzen:} ``Gleiche Richtung'', ``Tauziehen'', ``Rechtwinklig''
\item Ideal zur Demonstration vor der Klasse
\end{itemize}

\vspace{0.5cm}

\subsection*{6. Lösungen}

$\rightarrow$ Siehe separate Musterlösung (ML.pdf)

\vspace{0.3cm}

\textbf{Kurzübersicht Ergebnisse (18 BE gesamt):}

\begin{tabular}{ll}
Aufg. 1 (4P) & Kraftpfeile: 3N, 4N, 2{,}5N \\
Aufg. 2b) (4P) & $F_{res} = 7$ N (nach rechts) \\
Aufg. 3b) (5P) & $F_{res} = 2$ N (nach links, Team A gewinnt) \\
Aufg. 4a) (2P) & $F_{res} = 0$ N (Gleichgewicht) \\
Aufg. 5 (3P) & $F_{res} = 5$ N (Pythagoras: $\sqrt{25}$) \\
\end{tabular}

\vspace{0.5cm}

\hrule
\vspace{0.3cm}
\textit{Erstellt für  -- Physik Klasse 7}

\end{document}
